\chapter{Estructura del Código Fuente y Módulos}
\label{ch:apendice_codigo}
\renewcommand{\thesubsection}{\thesection.\arabic{subsection}}

En este apéndice se detalla la organización modular del código fuente del proyecto \texttt{neurolab}, la jerarquía de paquetes en Python 3.14+ y la cobertura de la suite de pruebas unitarias.

% =====================================================================
\section{Árbol de directorios del paquete \texttt{neurolab}}
\label{sec:arbol_codigo}

La estructura del código principal dentro de \texttt{neuromorphic\_lab/neurolab/} se organiza de la siguiente manera:

\begin{verbatim}
neurolab/
|-- __init__.py
|-- core/
|   |-- __init__.py
|   |-- base_device.py       # Clase base abstracta MemristorBase
|   +-- constants.py         # Constantes físicas universalmente usadas
|-- devices/
|   |-- __init__.py
|   |-- strukov.py           # Modelo lineal Strukov 2008
|   |-- yakopcic.py          # Modelo asimétrico Yakopcic 2011
|   +-- volatile.py          # Modelo volátil HfO2 con relajación
|-- neurons/
|   |-- __init__.py
|   |-- lif.py               # Neurona Leaky Integrate-and-Fire
|   +-- base_neuron.py       # Interfaz base de neuronas
|-- synapses/
|   |-- __init__.py
|   |-- memristive_synapse.py# Conexión sináptica basada en memristor
|   +-- plasticity.py        # Módulo de reglas STDP y R-STDP
|-- circuits/
|   |-- __init__.py
|   +-- hybrid_circuit.py    # Circuito híbrido Memristor-LIF
|-- crossbar/
|   |-- __init__.py
|   |-- core.py              # Clase Crossbar unificada (M x N)
|   |-- drivers.py           # Decodificadores y drivers de línea
|   +-- selector.py          # Esquema de selección V/2
|-- gui/
|   |-- __init__.py
|   |-- main_window.py       # Ventana principal PySide6
|   |-- base_view.py         # Canvas base genérico BaseCrossbarCanvas
|   |-- crossbar_view.py     # Panel multitab Crossbar
|   |-- crossbar_2x2_view.py # Vista interactiva 2x2
|   |-- crossbar_4x4_view.py # Vista interactiva 4x4
|   +-- widgets/             # Componentes Matplotlib y sliders
|-- io/
|   |-- __init__.py
|   |-- config_loader.py     # Carga/guardado de YAML/JSON
|   +-- exporter.py          # Exportador de CSV y figuras
+-- validation/
    |-- __init__.py
    +-- metrics.py           # Cálculo de MAE, RMSE, RelErr y R^2
\end{verbatim}

% =====================================================================
\section{Descripción de los módulos principales}
\label{sec:descripcion_modulos}

\begin{table}[H]
\caption{Resumen de módulos del paquete \texttt{neurolab}.}
\label{tab:modulos_neurolab}
\centering
\begin{tabular}{l p{10cm}}
\toprule
Módulo & Responsabilidad Principal \\
\midrule
\texttt{neurolab.core} & Define clases base abstractas, interfaces y constantes físicas universalmente usadas. \\
\texttt{neurolab.devices} & Implementa los modelos memristivos (Strukov, Yakopcic, Memristor Volátil). \\
\texttt{neurolab.neurons} & Implementa la neurona Leaky Integrate-and-Fire (LIF) y modelos adaptativos. \\
\texttt{neurolab.synapses} & Implementa la sinapsis memristiva y algoritmos de plasticidad (LTP/LTD, STDP, R-STDP). \\
\texttt{neurolab.circuits} & Ensambla el circuito híbrido sensor-neurona acoplando la corriente inyectada. \\
\texttt{neurolab.crossbar} & Gestiona matrices memristivas de dimensión $M \times N$, multiplicación VMM y esquema $V/2$. \\
\texttt{neurolab.gui} & Proporciona la interfaz gráfica PySide6 con renderizado dinámico interactivo. \\
\texttt{neurolab.io} & Gestiona el almacenamiento de datos, exportación en CSV y configuraciones YAML. \\
\texttt{neurolab.validation} & Contiene las funciones matemáticas para el cálculo de métricas ($R^2$, MAE, RMSE). \\
\bottomrule
\end{tabular}
\end{table}

% =====================================================================
\section{Resumen de la suite de pruebas unitarias (Pytest)}
\label{sec:suite_tests}

La calidad del código y la estabilidad de las refactorizaciones se verifican mediante una suite de 87 pruebas automatizadas organizadas por módulo:

\begin{table}[H]
\caption{Distribución de las 87 pruebas unitarias e integradas (todas pasando satisfactoriamente).}
\label{tab:resumen_tests}
\centering
\begin{tabular}{l c c}
\toprule
Directorio de Pruebas & Número de Pruebas & Estado \\
\midrule
\texttt{tests/core/} & 12 & PASS \\
\texttt{tests/devices/} & 18 & PASS \\
\texttt{tests/neurons/} & 10 & PASS \\
\texttt{tests/synapses/} & 14 & PASS \\
\texttt{tests/crossbar/} & 16 & PASS \\
\texttt{tests/gui/} & 8 & PASS \\
\texttt{tests/validation/} & 9 & PASS \\
\midrule
\textbf{Total} & \textbf{87} & \textbf{100\% Exitoso} \\
\bottomrule
\end{tabular}
\end{table}

% =====================================================================
\section{Fragmentos de Código Centrales}
\label{sec:fragmentos_codigo}

A continuación se exponen fragmentos representativos del núcleo matemático del simulador \texttt{neurolab}, ilustrando la aplicación de tipado estático, modularidad y las leyes físicas descritas en el marco teórico.

\subsection{Modelo Matemático de Strukov (\texttt{strukov.py})}
Implementación de la ecuación diferencial gobernante (dx/dt) encapsulada como un módulo matemático independiente.

\begin{verbatim}
class StrukovMathModel(BaseMathModel):
    """
    Modelo Matemático Determinista de Strukov et al. (Nature 2008).
    Ecuación diferencial de deriva iónica normalizada:
        dx/dt = (mu_v * R_on / D^2) * I(t)
    """
    def compute_dxdt(
        self, state: float, voltage: float, current: float, 
        electrical: ElectricalConfig, model_config: Optional[StrukovConfig] = None
    ) -> float:
        cfg = model_config if model_config is not None else StrukovConfig()
        
        # Computación de la tasa de cambio de la variable de estado
        return float((cfg.mu_v * electrical.r_on / (cfg.D ** 2)) * current)
\end{verbatim}

\subsection{Motor de Integración Base (\texttt{base\_device.py})}
El núcleo de integración numérica de Euler donde convergen el modelo matemático, las funciones ventana y la inyección de procesos estocásticos.

\begin{verbatim}
def step(self, dt: float, voltage: float) -> None:
    # 1. Resolver corriente (Ley de Ohm analógica con resistencia actual)
    current = voltage / self.R
    
    # 2. Computar derivada pura (dx/dt) desde el modelo matemático subyacente
    dx_dt = self.math_model.compute_dxdt(self.x, voltage, current, 
                                         self.electrical, self.model_config)
    
    # 3. Aplicar funciones ventana (ej. Biolek) y ruido termodinámico
    for modifier in self.modifiers:
        dx_dt = modifier.apply(self.x, voltage, current, dx_dt, dt, self)
        
    # 4. Integración explícita hacia adelante
    self.x += dx_dt * dt
    
    # 5. Saturación geométrica estricta a límites físicos
    if self.clip_x:
        self.x = max(0.0, min(1.0, self.x))
        
    self._history_x.append(self.x)
    self._history_r.append(self.R)
\end{verbatim}
